\documentclass[11pt,reqno]{amsart}
\usepackage[margin=1in]{geometry}
\usepackage[expansion=false]{microtype}
\usepackage{amssymb,mathtools,xcolor}
\usepackage[colorlinks=true,linkcolor=blue!50!black,citecolor=blue!50!black,urlcolor=blue!50!black]{hyperref}

\newtheorem{theorem}{Theorem}
\newtheorem{proposition}[theorem]{Proposition}
\newtheorem{lemma}[theorem]{Lemma}
\newcommand{\R}{\mathbb R}
\newcommand{\E}{\mathbb E}
\newcommand{\one}{\mathbf 1}
\newcommand{\ga}{\gamma}
\newcommand{\Phib}{\overline{\Phi}}

\hypersetup{
 pdftitle={Asymptotically Optimal Partitions for Agreement of Many Gaussian Samples},
 pdfauthor={Roy van Rijn}
}
\title[Agreement of many Gaussian samples]{Asymptotically Optimal Partitions for\\Agreement of Many Gaussian Samples}
\author{Roy van Rijn}
\subjclass[2020]{60E15, 60G15, 60F05}
\keywords{Gaussian noise stability, equicorrelated agreement, regular simplex, quantile slabs, low influences}
\date{}

\begin{document}
\begin{abstract}
We study the probability that $m$ equicorrelated Gaussian observations
all fall in the same cell of a partition into $k$ equal-measure sets.
For fixed $k\ge3$ and correlation $0<\rho<1$, the largest agreement
probability is asymptotic to $2H_m$ as $m\to\infty$, uniformly in the
ambient dimension. Here $H_m$ is the probability that all observations
fall in a single halfspace of mass $1/k$, and
$H_m=m^{-(1-\rho)/\rho+o(1)}$. Dividing one coordinate into $k$ equal-mass
intervals attains this bound asymptotically. The regular simplex has
agreement exponent $2(k-1)/k$ times the halfspace exponent. This gives
a counterexample to the many-observation simplex proposal of
Isaksson and Mossel and yields exactly balanced low-influence voting
rules with greater unanimous agreement than plurality for sufficiently
many correlated tallies.
\end{abstract}
\maketitle

\section{The agreement problem}

A partition of Gaussian space labels each observation by its cell.
Given many observations with the same positive pairwise correlation,
we ask for the largest probability that they all receive one label.

Fix $k\ge3$, $0<\rho<1$, and $d\ge1$. Write $\ga_d$ for standard Gaussian
measure on $\R^d$. A partition $A=(A_1,\ldots,A_k)$ is \emph{balanced}
if $\ga_d(A_i)=1/k$ for each $i$. For jointly Gaussian observations
$Y_1,\ldots,Y_m$ with mean zero, covariance $I_d$ per observation, and
$\operatorname{Cov}(Y_j,Y_\ell)=\rho I_d$ for $j\ne\ell$, set
\[
 M_m(A)=\sum_{i=1}^k\Pr(Y_1,\ldots,Y_m\in A_i).
\]
Let $H_m$ be the corresponding probability for a single halfspace of
mass $1/k$. Let $Q$ be the partition into the inverse images of the $k$
quantile intervals of a standard normal under the first coordinate.
The two end cells of $Q$ are opposite halfspaces.

\begin{theorem}\label{thm:agreement}
For fixed $k$ and $\rho$ as above, in every $d\ge1$,
\[
 2H_m\ \le\ M_m(Q)\ \le\ \sup_{A\ \mathrm{balanced}} M_m(A)
       \ =\ (2+o(1))H_m.
\]
Here $H_m=m^{-\alpha+o(1)}$, with $\alpha=(1-\rho)/\rho$, as $m\to\infty$.
The error in the supremum is uniform over $d\ge1$.
\end{theorem}

When $m=2$, this is the agreement objective in the Standard Simplex
Conjecture of Isaksson and Mossel \cite[Conjecture~1.4]{IM}. Their
Section~8 also proposes a many-observation extension. In the
equicorrelated setting of that extension, the simplex can be beaten for
every fixed $k\ge3$ and $\rho\in(0,1)$ when $m$ is sufficiently large.
For $d\ge k-1$, let $T$ denote the partition whose cells are the winning
regions of a centered regular simplex, extended cylindrically to $\R^d$.

\begin{proposition}\label{prop:simplex}
With $k$ and $\rho$ fixed and $d\ge k-1$, the centered regular-simplex
partition satisfies
\[
 M_m(T)=m^{-\beta\alpha+o(1)},\qquad
 \beta=\frac{2(k-1)}{k}.
\]
\end{proposition}

Together the two estimates give $M_m(Q)/M_m(T)=m^{(\beta-1)\alpha+o(1)}$.
For example, when $k=3$ and $\rho=1/2$, the halfspace exponent is $1$
and the simplex exponent is $4/3$. The exponent gap comes from the
need to satisfy several facet inequalities to make the conditional
probability of a simplex cell close to one. In the upper bound for an
arbitrary partition, a different geometric effect forces the constant
$2$. Conditional probabilities close to one for different labels
require well-separated common Gaussian means. Three such sets cannot
all have large Gaussian measure. The single-cell comparison of
Isaksson and Mossel \cite[Theorem~1.2]{IM} then controls the two
remaining terms. The result gives the asymptotic value of the
supremum. Identifying exact optimizers at finite $m$ and classifying
the families that attain the asymptotic value are separate questions.

\section{Halfspace and simplex posterior probabilities}

Write $s=\sqrt{\rho}$, $\sigma=\sqrt{1-\rho}$, and $a=\sigma/s$.
One can represent the observations as $Y_j=sX+\sigma Z_j$, with
$X,Z_1,\ldots,Z_m$ independent and standard Gaussian in $\R^d$.
For each cell set
\[
 p_i(x)=\Pr(sx+\sigma Z_1\in A_i),\qquad
 M_m(A)=\sum_{i=1}^k \E[p_i(X)^m].
\]
In the displayed identity, the observations are conditionally independent
given $X$. Gaussian smoothing of a cell indicator gives a continuous
function $p_i$. We may choose Borel cell representatives since changing
a cell on a Gaussian null set leaves every conditional probability
unchanged.

Write $\Phi$ for the standard normal distribution function and
$\Phib=1-\Phi$. For the halfspace $\{y:y_1>h\}$ of mass $1/k$,
where $h=\Phi^{-1}(1-1/k)$, its conditional probability is
$p_H(x)=\Phi((sx_1-h)/\sigma)$. For $\epsilon\in(0,1/2)$ put
$t=\Phi^{-1}(1-\epsilon)$. Then
\begin{equation}\label{eq:halfspace-tail}
 \Pr(p_H(X)>1-\epsilon)=\Phib(at+h/s).
\end{equation}
The Gaussian tail estimate $\log\Phib(t)=-t^2/2+o(t^2)$ gives
$\Pr(p_H(X)>1-\epsilon)=\epsilon^{\alpha+o(1)}$ as
$\epsilon\downarrow0$.

\begin{lemma}\label{lem:posterior-powers}
If $P$ takes values in $[0,1]$ and
$\Pr(P>1-\epsilon)=\epsilon^{\eta+o(1)}$ as $\epsilon\downarrow0$
for some $\eta>0$, then $\E[P^m]=m^{-\eta+o(1)}$.
\end{lemma}

\begin{proof}
On $\{P>1-1/m\}$, the factor $P^m$ is bounded below by
$(1-1/m)^m$ for $m\ge2$. For the reverse bound, split at
$1-C\log(m)/m$ with fixed $C>\eta$. Below this level the $m$th power
is at most $m^{-C}$. The probability of exceeding the level is
$m^{-\eta+o(1)}$ by the assumed tail.
\end{proof}

Since $H_m=\E[p_H(X)^m]$, Lemma~\ref{lem:posterior-powers} gives the
halfspace power law of Theorem~\ref{thm:agreement}.

\begin{proof}[Proof of Proposition~\ref{prop:simplex}]
Identify the active space with $E=\one^\perp\subset\R^k$ and write
$e_1,\ldots,e_k$ for the coordinate vectors of $\R^k$. The cell
with label $i$ is the intersection of the halfspaces with inward
unit normals $n_{ij}=(e_i-e_j)/\sqrt2$, $j\ne i$.
If $p_i(x)>1-\epsilon$, each containing halfspace has conditional
probability greater than $1-\epsilon$. Thus
$x\cdot n_{ij}>at$ for all $j\ne i$.
The sum of the $k-1$ normals has length $\sqrt{k(k-1)/2}$.
If $u_i$ is its unit direction, summing the inequalities gives
\begin{equation}\label{eq:simplex-upper}
 x\cdot u_i>\sqrt{\beta}\,at,\qquad
 \Pr(p_i(X)>1-\epsilon)\le\Phib(\sqrt{\beta}\,at).
\end{equation}

For the converse let $t'=\Phi^{-1}(1-\epsilon/(k-1))$. Each
$n_{ij}$ has inner product $1/\sqrt{\beta}$ with $u_i$.
On the event
\[
 X\cdot u_i>\sqrt{\beta}(at'+1),\qquad
 \big|X_E-(X\cdot u_i)u_i\big|\le1,
\]
where $X_E$ denotes the projection onto $E$, every margin exceeds
$at'$. The union bound over the $k-1$ halfspaces gives
$p_i(X)>1-\epsilon$. The two conditions are independent, and
the second has a fixed positive probability in the $k-2$ dimensional
space perpendicular to $u_i$ within $E$. The Gaussian tail estimate
gives $t'/t\to1$. With \eqref{eq:simplex-upper} we conclude that
\[
 \Pr(p_i(X)>1-\epsilon)=\epsilon^{\beta\alpha+o(1)}
 \qquad(\epsilon\downarrow0).
\]
Lemma~\ref{lem:posterior-powers} gives the exponent for one cell, and
summing over the fixed $k$ preserves it.
\end{proof}

\section{Only two cells can contribute at the halfspace scale}

For a balanced partition and $0<\epsilon<1/2$ let
\[
 D_i=\{x:p_i(x)>1-\epsilon\},\qquad t=\Phi^{-1}(1-\epsilon).
\]
We first bound the measure of the third largest of the $D_i$.
The geometry of the conditional laws supplies a bound uniform
over all ambient dimensions.

\begin{lemma}\label{lem:separation}
For distinct labels $i,j$, all $x\in D_i$ and $y\in D_j$ satisfy
$|x-y|\ge2at$. For three distinct labels $i,j,\ell$,
\[
 \big(\ga_d(D_i)\ga_d(D_j)\ga_d(D_\ell)\big)^{1/3}
 \le \exp(-2a^2t^2/3).
\]
\end{lemma}

\begin{proof}
Under the Gaussian law with mean $sx$ and covariance $\sigma^2 I_d$,
the set $A_i$ has measure greater than $1-\epsilon$. Under the law
with mean $sy$ and the same covariance it has measure less than
$\epsilon$, since $A_j$ has measure greater than $1-\epsilon$.
The test using $A_i$ therefore has sum of errors less than $2\epsilon$.
The smallest possible sum of errors is the integral of the pointwise
minimum of the two densities. For $x\ne y$ the two densities cross on
the plane through the midpoint of their means, perpendicular to
$x-y$, so projection gives the minimum as
$2\Phib(s|x-y|/(2\sigma))$. The same formula gives one when $x=y$.
It follows that $s|x-y|/(2\sigma)\ge t$.

Now take $x_1,x_2,x_3$ in the three respective sets and put
$z=(x_1+x_2+x_3)/3$. The separation just proved gives
\[
 \frac13\sum_{\nu=1}^3 |x_\nu|^2-|z|^2
 =\frac19\sum_{\mu<\nu}|x_\mu-x_\nu|^2
 \ge \frac{(2at)^2}{3}.
\]
Write $\varphi_d$ for the standard Gaussian density. The geometric
mean of the three restricted densities
$\varphi_d 1_{D_i}$, $\varphi_d 1_{D_j}$, $\varphi_d 1_{D_\ell}$,
evaluated at their respective points, is bounded by
$\exp(-2a^2t^2/3)\varphi_d(z)$. If any point lies outside its
assigned set, the geometric mean is zero and the same bound holds.
The three-function Pr\'ekopa--Leindler inequality
\cite[Theorem~3]{Prekopa} applied with weights $1/3$ gives the
asserted product bound.
\end{proof}

\begin{proof}[Proof of Theorem~\ref{thm:agreement}]
By Lemma~\ref{lem:separation}, at most two labels satisfy
$\ga_d(D_i)>\exp(-2a^2t^2/3)$. Choose two labels covering all
those that do. The Exchangeable Gaussians Theorem of Isaksson and
Mossel \cite[Theorem~1.2]{IM} bounds the entire term
$\Pr(Y_1,\ldots,Y_m\in A_i)$ by $H_m$ for each of the two chosen
labels. Its hypotheses apply to our jointly Gaussian observations
with common positive cross covariance and to the repeated set $A_i$
of mass $1/k$. The identical-cell case also follows from
Borell's earlier Gaussian rearrangement theorem
\cite[Theorem~3.1]{Borell}.

For every other label, integrate the $m$th power of $p_i$ separately
on $D_i$ and its complement. This gives the dimension-independent
finite-$m$ bound
\begin{equation}\label{eq:upper}
 M_m(A)\ \le\
 2H_m+(k-2)\big((1-\epsilon)^m+\exp(-2a^2t^2/3)\big).
\end{equation}
Choose $\epsilon=C\log(m)/m$ for any fixed $C>\alpha$. The first
term inside the parentheses is at most $m^{-C}$. Since
$t^2=2\log(1/\epsilon)+o(\log(1/\epsilon))$, the second term is
$m^{-4\alpha/3+o(1)}$. Both are $o(H_m)$ by the halfspace
estimate \eqref{eq:halfspace-tail} and Lemma~\ref{lem:posterior-powers}. The two end cells of the quantile partition
each contribute $H_m$, so $M_m(Q)\ge2H_m$. This gives both the
asymptotic supremum and uniformity in dimension.
\end{proof}

\section{Exactly balanced voting rules}

Consider $n$ independent voters. For each voter, draw a common ballot
uniformly from $\{1,\ldots,k\}$. For each of the $m$ copies
independently, keep the common ballot with probability $s=\sqrt\rho$.
Otherwise draw a new uniform ballot independently. The ballot profile
in each copy is uniform. For a rule
$f:\{1,\ldots,k\}^n\to\{1,\ldots,k\}$, its influence at a voter is
the probability that independently resampling that vote in a uniform
profile changes $f$. Exact balance means every label is chosen with
probability $1/k$ on a uniform profile.

\begin{theorem}\label{thm:voting}
Fix $k\ge3$ and $\rho\in(0,1)$. For every sufficiently large fixed
$m$ there are exactly balanced deterministic rules $f_n$ whose maximum
voter influence tends to zero as $n\to\infty$. Their limiting
probability of unanimous agreement across the $m$ copies is strictly
greater than for plurality in the same model.
\end{theorem}

\begin{proof}
Write $b=1/k$ and for each copy center the count of votes for
candidate $1$ at $nb$ and divide by $\sqrt{nb(1-b)}$. For a single
voter the indicator of a vote for candidate $1$ has variance
$b(1-b)$. Its conditional mean given the common ballot $B$ is
$s1_{\{B=1\}}+(1-s)b$. Distinct copies are conditionally independent
given $B$, so the cross covariance is $s^2 b(1-b)=\rho b(1-b)$.
The multivariate central limit theorem, at fixed $m$, gives the
one-dimensional equicorrelated Gaussian model for the normalized
counts across copies.

Apply the $k$ Gaussian quantile intervals to these counts to get a
rule $g_n$. By convergence at the thresholds, its label masses tend
to $1/k$ and its all-copy agreement tends to $M_m(Q)$. A resampled
vote moves the normalized count by at most $(nb(1-b))^{-1/2}$.
It changes the label only within that distance of a threshold.
For each fixed positive neighborhood size the limit probability
is its Gaussian measure. Letting that size decrease to zero
shows that the maximum voter influence tends to zero.

To get exact balance, change labels on a minimal set of profiles
by assigning surplus profiles to deficient labels. The target
number $k^{n-1}$ of profiles per label is integral. The fraction
$\delta_n$ of profiles changed equals half the sum of the
absolute discrepancies of the label masses from $1/k$ and tends
to zero. For a given voter, the original profile and the
independently resampled profile both have uniform marginals. The
influence increases by at most $2\delta_n$. For $m$ copies the
probability of any change from $g_n$ is at most $m\delta_n$, so
the agreement limit of the exactly balanced rule $f_n$ remains
$M_m(Q)$.

For plurality use the centered count vector in each copy. Its
one-voter covariance is $k^{-1}$ times the orthogonal projection
onto $\one^\perp$, and the cross-copy covariance is $\rho$
times the same matrix, by conditional independence as above.
After normalization, the multivariate central limit theorem
gives $m$ standard Gaussians on $\one^\perp$ with cross covariance
$\rho I$. The largest coordinate is the simplex label and the
Gaussian tie planes have measure zero. Thus plurality agreement
converges to $M_m(T)$ for any rule for breaking ties. Plurality
itself can be balanced at every $n$ by selecting the candidate of the
first voter whose ballot names a tied winner. This choice commutes
with permutations of the candidates.
By Theorem~\ref{thm:agreement} and Proposition~\ref{prop:simplex},
$M_m(Q)>M_m(T)$ for sufficiently large $m$. For each such fixed $m$, the two limits in $n$ give the
strict comparison.
\end{proof}

\paragraph{AI assistance and review.}
ChatGPT (OpenAI) assisted with mathematical arguments, literature checks,
exposition, proof review and verification code.

\begin{thebibliography}{9}
\bibitem{Borell}
C.~Borell, \emph{Geometric bounds on the Ornstein--Uhlenbeck velocity
process}, Z. Wahrscheinlichkeitstheorie verw. Gebiete \textbf{70}
(1985), 1--13.
\href{https://doi.org/10.1007/BF00532234}{doi:10.1007/BF00532234}.
\bibitem{IM}
M.~Isaksson and E.~Mossel,
\emph{Maximally stable Gaussian partitions with discrete applications},
Israel J. Math. \textbf{189} (2012), 347--396.
\href{https://doi.org/10.1007/s11856-011-0181-7}{doi:10.1007/s11856-011-0181-7}.
See also \href{https://arxiv.org/abs/0903.3362v3}{arXiv:0903.3362v3}.
\bibitem{Prekopa}
A.~Pr\'ekopa, \emph{On logarithmic concave measures and functions},
Acta Sci. Math. (Szeged) \textbf{34} (1973), 335--343.
\href{https://acta.bibl.u-szeged.hu/14411/}{Journal archive}.
\end{thebibliography}
\end{document}

